\documentclass[12pt]{article}
\usepackage[margin=1in]{geometry}
\usepackage{amsmath}
\usepackage{natbib}
\usepackage{booktabs}
\usepackage{longtable}
\usepackage{hyperref}

\title{Underwriters and Wholesalers of United States Federal Debt,\\
1775--1940 (and After):\\
Who Marketed the Government's Bonds, and How Much Went to Europe}
\author{George J. Hall and Thomas J. Sargent\thanks{Background memorandum
prepared with Claude. Quantitative claims marked ``verified'' in the text
were checked against the primary and near-primary sources flagged
\texttt{[PRIMARY]} in the accompanying \texttt{bond\_underwriters.bib};
page numbers for several standard monographs were not confirmed against the
physical volumes and are flagged for checking before print.}}
\date{\today}

\begin{document}

\maketitle

\begin{abstract}
This memorandum surveys the institutions and individuals who acted as
\emph{underwriters}, \emph{wholesalers}, or \emph{marketers} of initial
offerings of United States federal debt from the Revolution through 1940,
and it assembles what is known about the fraction of that debt placed or held
in Europe. The story divides into five regimes. (1) In the Revolution and
Confederation (1775--1790), thirteen state loan offices, the Continental loan
office, Robert Morris and his Bank of North America, and---decisively for
hard money---the French crown and a consortium of Amsterdam bankers
(Willink; Van Staphorst; De la Lande \& Fynje) placed the debt. (2) From 1790
to 1836 the two Banks of the United States served as fiscal agents while Dutch
and then British investors absorbed a remarkable share of the funded debt---a
share that peaked near $53$--$56$ percent in $1803$, the year of the
Louisiana Purchase. (3) From the demise of the Second Bank (1836) to the Civil
War, private houses---above all Corcoran \& Riggs, with London houses such as
Baring Brothers as a secondary outlet---wholesaled the government's loans. (4) Jay Cooke \& Co.
pioneered mass retail distribution of the Civil War debt (1862--1865); after
Cooke's failure in the Panic of 1873, leadership passed to Drexel, Morgan \&
Co. and its London parent J.\,S. Morgan \& Co., which dominated the refunding
syndicates of the 1870s and the gold-reserve rescue of 1895. (5) The First
World War reversed the historic pattern: the Liberty Loans were sold directly
to tens of millions of American citizens through the Federal Reserve Banks,
while J.\,P. Morgan \& Co. acted as purchasing agent for Britain and France
and organized the 1915 Anglo-French loan---America had become the world's
creditor. We close with a consolidated table of the European-held fraction of
the federal debt across a century and a half.
\end{abstract}

\section{Introduction}

A government that borrows must find someone to place its paper. Over the
century and a half between the Revolution and the Second World War the United
States relied on a changing cast of intermediaries: colonial-style loan
offices, a national bank acting as fiscal agent, foreign sovereigns and
foreign bankers, private ``wholesalers'' who bought whole loans at a discount
and retailed them onward, syndicates of investment bankers, and---finally---the
Federal Reserve System selling directly to the public. This memorandum
identifies those intermediaries regime by regime and, wherever the record
allows, states the fraction of each era's debt that came to rest in Europe.

Two questions organize the account. \emph{Who underwrote or wholesaled the
initial offerings?} And \emph{what share was sold or held abroad?} The two are
connected: for most of the nineteenth century the United States was a
capital-importing, ``emerging'' economy whose government could not have funded
its wars and its debt conversions without European savings, and the
intermediaries who mattered most were precisely those---the Amsterdam houses,
Baring Brothers, the Rothschilds through August Belmont, J.\,S. Morgan \& Co.
in London---with a foot on both sides of the Atlantic.

\section{The Revolution and the Confederation, 1775--1790}

\subsection{State and Continental loan offices}

Congress had no power to tax. It financed the Revolution first by emitting
Continental currency (the ``Continentals''), then, as that currency
depreciated toward worthlessness, by \emph{borrowing}. In 1776 Congress
established a \emph{Continental loan office} in each state, staffed by a
Continental loan officer, to sell ``loan-office certificates''---interest-bearing
IOUs of the Confederation---to citizens \citep{ferguson1961power}. The state
loan offices were the retail distribution network of the Revolution: they were
the counters at which patriotic (and speculative) citizens exchanged
depreciating money for the promise of future interest, and they doubled as the
registers of the domestic debt. Alongside the Continental certificates, the
individual states issued their own revolutionary debts, which Hamilton would
later propose the federal government assume.

\subsection{Robert Morris and the Bank of North America}

The near-collapse of Continental finance in 1780--1781 produced the
Confederation's one experiment in centralized fiscal management. Robert Morris,
appointed Superintendent of Finance in 1781, secured the incorporation of the
\emph{Bank of North America}---approved by Congress on 26 May 1781, chartered by
ordinance of 31 December 1781, and opened in Philadelphia on 7 January
1782---the first commercial bank in the United States, with a capital of
\$400,000 in 1{,}000 shares of \$400 \citep{rappleye2010robert}. Tellingly, most
of that capital came not from private subscribers but from the government
itself: Morris subscribed the specie proceeds of the French and Dutch loans, so
that ``the proceeds of a foreign loan\,\ldots\ did become the source of most of
[the Bank's] capital'' \citep[9]{sylla2011financial}. The Bank lent to the
government, gave its short-term paper (``Morris notes'') currency, and steadied
public finance; it was in that narrow sense the first institution to
``wholesale'' the credit of the Union, though its resources were small and its
life as a national fiscal agent brief \citep{ferguson1961power}. Morris leaned
on private brokers---most famously Haym Salomon, broker to the Office of Finance,
who sold bills of exchange on France and Holland and placed government
securities to convert the foreign subsidies into usable funds.

\subsection{The French, Spanish, and Dutch loans}

The hard money that actually paid for the Revolution came from Europe. The
whole amount received from France during the war, in loans and subsidies
together, was $45$ million livres, about \$8.17 million---of which roughly $34$
million livres were repayable \emph{loans} (the royal advances consolidated in
the contract of 16 July 1782 at $5$ percent, a French-guaranteed Holland loan of
$10$ million livres signed by Vergennes and Franklin at $4$ percent, and a $6$
million-livre loan of 1781) and $11$ million livres were outright \emph{gifts}
(a French subsidy of $10$ million plus a Spanish subsidy of $1$ million)
\citep[299--308]{bayley1881national}. The very first Continental loan, in 1777,
came from the \emph{Farmers-General} of France---a $1$ million-livre advance
against $5{,}000$ hogsheads of tobacco, arranged by Franklin and
Deane---while Beaumarchais's front company \emph{Roderigue, Hortalez \& Co.} was
the covert conduit for early aid \citep[304--305]{bayley1881national}. Spain,
petitioned by John Jay for a \$5 million loan, gave almost nothing beyond the
subsidy: the Spanish debt principal was just \$174{,}011
\citep[306]{bayley1881national}.

As independence approached, the decisive \emph{new} source was Amsterdam.
Beginning with John Adams's breakthrough loan of 1782, the United States
borrowed in the Dutch capital market through a consortium of three houses acting
together---\emph{Wilhem \& Jan Willink}, \emph{Nicolaas \& Jacob Van Staphorst},
and \emph{De la Lande \& Fynje} \citep{vanwinter1977american,
riley1980international}. That first loan was for $5{,}000{,}000$ guilders
(\$2 million), to run ten years at $5$ percent, its contract concluded 11 June
1782 and confirmed by Congress that September; it sold only gradually, the full
five million not taken up until 1786 \citep[311]{bayley1881national}. Adams
contracted four such loans (1782, 1784, 1787, 1788) totaling $9$ million
guilders; in all, \emph{eleven} United States loans were floated in Amsterdam
between 1782 and 1794 for roughly $29$--$30$ million guilders, and the
consortium served as ``the financial agents of the United States in Holland''
for some forty years \citep[311--317]{bayley1881national}. The houses' all-in
charges commonly ran about $4\tfrac{1}{2}$ percent, though they reached roughly
$8$ percent on the 1787 loan \citep[311--317]{bayley1881national}; on their own
account they also speculated heavily in the deeply discounted domestic debt in
anticipation of its funding \citep{syllawilsonwright2006}.

At the launch of the new government, Hamilton's First Report on Public Credit
(January 1790) reckoned the foreign-currency debt---owed chiefly to France and
to the Dutch bankers---at \$11{,}710{,}378, of a federal debt of
\$54{,}124{,}464; that foreign-origin debt was thus about a fifth ($22$ percent)
of the federal debt, or nearer a seventh once the roughly \$25 million of state
debts Hamilton proposed to assume are added to the base
\citep[12]{sylla2011financial}. This ``foreign-origin'' debt should not be
confused with foreign \emph{ownership} of the funded debt, which was a different
and, as we shall see, much larger magnitude by 1803.

\section{The Amsterdam Flotations in Detail, 1782--1794}
\label{sec:dutch}

Because the Amsterdam loans were the mechanism by which the young republic first
established, and then defended, its credit, they repay a closer look. Rafael
Bayley's Treasury monograph records eleven United States loans floated in
Amsterdam between 1782 and 1794, all through the same consortium---Willink; Van
Staphorst (from 1790 styled Van Staphorst \& Hubbard); and, until its failure in
1787, De la Lande \& Fynje. Table~\ref{tab:dutch} sets them out loan by loan,
converting guilders to dollars at Bayley's constant rate of $2.5$ florins to the
dollar \citep[310--321]{bayley1881national}.

The sequence has a clear internal logic. The first four loans (1782--1788),
contracted by John Adams---and, for the last, jointly with Jefferson at The
Hague---were the Confederation's lifeline, totaling $9{,}000{,}000$ guilders
(\$3.6 million). They sold slowly and dearly: the pioneering loan of 1782 took
four years to place, the 1784 loan had to be dressed up as a part-lottery to
attract subscribers, and by the 1787 loan---floated merely to pay the interest
falling due on its predecessors, with De la Lande \& Fynje now bankrupt and out
of the syndicate---the houses' all-in charge had risen to about $8$ percent
\citep[310--316]{bayley1881national}. The contrast after 1790 is striking. Once
Hamilton's funding system had restored the government's credit, William Short
could drive the bankers' commissions down toward $4$ percent and place paper with
an ease unimaginable a decade earlier: the loan of February 1791 was ``taken up
within two hours of publication---a celerity of which there had been no instance
before'' \citep[317--318]{bayley1881national}. The seven new-government loans
(1790--1794) raised $23{,}500{,}000$ guilders (\$9.4 million), the proceeds going
to buy up the domestic debt under the sinking fund, to pay off Spain and the
foreign-officer debt, and---loan by loan---to service and retire the earlier
French and Dutch obligations themselves. A parallel Antwerp loan of November 1791
(3 million Brabant florins at $4\tfrac{1}{2}$ percent, through C.\,J.\,M. De Wolf
\& Co.) was remitted directly to Paris to meet installments on the French debt;
Bayley excludes it from the Holland total \citep[318--321]{bayley1881national}.

\begin{table}[htbp]
\centering
\caption{United States loans floated in Amsterdam, 1782--1794 (guilders converted
at 2.5\,f.\ $=$ \$1)}
\label{tab:dutch}
\footnotesize
\begin{tabular}{l r r c c l p{3.4cm}}
\toprule
Contract & Florins & Dollars & Coupon & Charge & By & Purpose / note \\
\midrule
11 Jun 1782 & 5{,}000{,}000 & 2{,}000{,}000 & $5\%$ & $4.5\%$ & Adams & First U.S. loan in Holland; placed only by 1786 \\
9 Mar 1784  & 2{,}000{,}000 & 800{,}000   & $4\%^{*}$ & $4.5\%$ & Adams & Part-lottery; effective $\sim\!5\%$ \\
18 May 1787 & 1{,}000{,}000 & 400{,}000   & $5\%$ & $\sim\!8\%$ & Adams & Pays interest then due; De la Lande \& Fynje fails, exits \\
1788        & 1{,}000{,}000 & 400{,}000   & $5\%$ & --- & Adams & Bridges the 1788--90 interregnum (with Jefferson) \\
\midrule
\multicolumn{6}{l}{\emph{Confederation subtotal (Adams)}} & 9{,}000{,}000\,f. $=$ \$3.6\,M \\
\midrule
12 Nov 1790 & 3{,}000{,}000 & 1{,}200{,}000 & $5\%$ & $4.5\%$ & Hamilton & First loan of the new government \\
15 Feb 1791 & 2{,}500{,}000 & 1{,}000{,}000 & $5\%$ & $4\%$ & Short & ``Taken up within two hours'' \\
Sep 1791    & 6{,}000{,}000 & 2{,}400{,}000 & $5\%$ & $4\%$ & Short & Cap lifted to seize a good market \\
Dec 1791    & 3{,}000{,}000 & 1{,}200{,}000 & $4\%$ & $5.5\%$ & Short & Lower coupon, higher commission \\
21 Mar 1792 & 2{,}950{,}000 & 1{,}180{,}000 & $4\%$ & $5\%$+ & Short & Closed short of 3\,million \\
1793        & 5{,}000{,}000 & 2{,}000{,}000 & $5\%$ & $3.5\%$ & bankers & \emph{Reloan} to meet the 1782 installment \\
10 Apr 1794 & 3{,}000{,}000 & 1{,}200{,}000 & $5\%$ & $4\%$ & Willink/V.S. & ``Closes the account of the Holland loans'' \\
\midrule
\multicolumn{6}{l}{\emph{New-government subtotal (Hamilton/Short)}} & 23{,}500{,}000\,f. $=$ \$9.4\,M \\
\bottomrule
\end{tabular}

\smallskip
\footnotesize Source: \citet[310--321]{bayley1881national}. $^{*}$Nominal $4\%$
plus lottery premiums, which Bayley reckons at an effective $\sim\!5\%$. Bayley's
stated grand total is $32{,}500{,}000$\,f.\ ($\$13$ million), of which
$9{,}000{,}000$\,f.\ were Confederation loans and $23{,}500{,}000$\,f.\ loans of
the present government.
\end{table}

How many loans, and how much in all? The answer depends on how one treats the
loan of 1793, which Bayley himself labels a ``reloan or continuation'' rather than
fresh borrowing---money raised only to meet the first installment coming due on
the original 1782 loan. Counting it, there were eleven loans; Bayley's round total
is \$13 million ($32.5$ million guilders), though his own itemized principals
actually sum to about $34.5$ million guilders, so the round figure is best read as
an editorial approximation. Excluding the 1793 refinancing, the ``new-money''
total falls to roughly $29.5$ million guilders---which is why the standard
scholarly source \citep{vanwinter1977american} gives about $29$ million and the
Amsterdam exchange historians about $30.5$ million. The apparent conflict among
``$29$,'' ``$30.5$,'' and ``$32.5$ million guilders'' thus dissolves once the 1793
continuation is identified \citep[310--321]{bayley1881national}.

\section{Funding, Assumption, and the Bank of the United States, 1790--1811}

\subsection{Hamilton's funding and the First Bank as fiscal agent}

Hamilton's Funding Act of 1790 refunded the Confederation's domestic and
assumed state debts into three new classes of federal stock (the ``sixes,''
the ``deferred sixes,'' and the ``threes''), while the foreign debt was
serviced and progressively refinanced through new Amsterdam loans. The
principal domestic institution supporting the market was the \emph{First Bank
of the United States} (chartered 1791, capital \$10 million, charter expiring
1811). The First Bank was the government's fiscal agent---holder of Treasury
balances, transfer agent, and a large and steadying holder and lender against
federal stock---rather than an ``underwriter'' in the later syndicate sense
\citep{hammond1957banks, sylla2011financial}. Domestic financiers such as
William Bingham of Philadelphia dealt heavily in the new stock, but the marginal
buyer who set the price was increasingly European.

\subsection{The European high-water mark, 1803}

The share of the federal debt held abroad rose steeply through the 1790s and
peaked in 1803. According to the reconciliation of \citet{syllawilsonwright2006},
built from Blodget's 1806 estimates and Bayley's Treasury figures, foreigners
held about \$43.4 million of the \$81.4 million federal debt in 1803---roughly
\textbf{53 percent} (Mira Wilkins's headline figure is \textbf{56 percent})
\citep{wilkins1989history, blodget1806economica}. The composition had shifted
from Dutch to British predominance, with the House of Baring now prominent.
Two facts dramatize the dependence on European capital in this year:
foreigners held about $62$ percent of the stock of the (first) Bank of the
United States, and the \$11.25 million ``Louisiana sixes'' issued to pay France
for the Louisiana Purchase were, at issue, taken up essentially \emph{entirely}
abroad---the Barings and Hope \& Co. of Amsterdam underwrote the Louisiana
stock and placed it in Europe \citep{syllawilsonwright2006}. The United States
had, in effect, borrowed in Europe to buy a continent from a European power.

\section{War, a Second Bank, and the Extinction of the Debt, 1812--1836}

\subsection{The war loans of 1812--1815 and the Girard--Astor--Parish syndicate}

The War of 1812 was financed, badly, by loans that the Treasury struggled to
sell as the war went against it and New England capital abstained. The pivotal
episode---and the first true \emph{private underwriting syndicate} in U.S.
history---was the \$16 million loan of 1813 (authorized by the act of 8 February
1813). Public subscription reached only about \$6 million, leaving the bulk to
three financiers: \emph{Stephen Girard} of Philadelphia, \emph{John Jacob Astor}
of New York, and \emph{David Parish}. Girard and Parish took \$7{,}055{,}800 on
joint account and Astor and his associates about \$2{,}056{,}000, at a price of
$88$---a $12$ percent discount on the $6$ percent stock, which, as Parish noted
in his own correspondence, ``yields an interest of $7\tfrac{1}{2}$ pct''
\citep{porter1931astor, adams1978finance, dewey1903financial}. The bet paid off
quickly: the stock soon traded at a premium of five or six points, ``promising a
handsome profit'' on top of the coupon. Girard's own bank (successor to the
First Bank's Philadelphia premises) and Astor's mercantile fortune were the
balance-sheet capacity behind the deal, which kept the government solvent at its
lowest ebb.

\subsection{The Second Bank of the United States, 1816--1836}

The Second Bank of the United States (chartered 1816, capital \$35 million)
restored a national fiscal agent. Like the First Bank, it served principally as
the depository of federal funds, the transfer and interest-paying agent for the
public debt, and a stabilizer of the market for government stock rather than as
a syndicate underwriter \citep{catterall1903second}. Its era, however, was also
the era of \emph{debt extinction}: under sustained surpluses the national debt
was paid down to essentially zero by 1835---the only time in American history
the federal debt has been extinguished \citep{lane2014nation}. With no debt to
market and, after 1836, no national bank, the institutional apparatus of
federal borrowing had to be rebuilt from private materials when deficits
returned. By 1818 the foreign-held share of the (now small) federal debt had
fallen to about $26$ percent \citep{wilkins1989history}.

\section{The Private Wholesalers of the Antebellum Treasury, 1836--1861}

\subsection{Selling loans without a national bank}

With the Second Bank gone and the Independent Treasury not yet permanent, the
government financed the recurrent deficits of 1837--1843 (expenditure exceeded
revenue by \$31.3 million over the four years to 1841) by two expedients: rolling
over short-term \emph{Treasury notes}, and selling \emph{funded loans} by public
advertisement and invited competitive proposals---no single wholesaler yet
dominated. The loan of 1841 (\$12 million authorized at $\le 6$ percent, no stock
below par) ``does not appear to have met with much favor,'' and only
\$5{,}672{,}077 could be placed \citep[363]{bayley1881national}. Only after the
act of 1842 permitted sale \emph{below} par did the market clear: \$8{,}343{,}886
of stock was sold for \$8{,}301{,}468 in cash, a shade under par
\citep[363--364]{bayley1881national}. It was in exactly this environment that
W.\,W. Corcoran got his start, taking a \$5 million loan from Secretary Forward.

\subsection{Corcoran \& Riggs and the Mexican War loans}

The Mexican War of 1846--1848 made \emph{Corcoran \& Riggs} of Washington
(W.\,W. Corcoran and George Washington Riggs) the pre-eminent wholesaler of the
federal debt. Of the loan of 1847 (act of 28 January 1847; \$28.2 million of $6$
percent stock ultimately issued, most sold at a small premium
\citep{bayley1881national}), Corcoran, singly or in association, cornered about
\$14.7 million of the roughly \$16.5 million public offering---$80$ to $90$
percent---and the firm ``earned more than \$250{,}000'' on it
\citep[45--49]{cohen1971business}. Of the loan of 1848 (act of 31 March 1848;
\$16 million at $6$ percent) the firm again took about \$14 million
\citep[49--51]{cohen1971business}. Philip Hone marveled
that Corcoran \& Riggs had ``made a princely fortune by taking the whole of the
last government six per cent loan.''

Two corrections to the conventional telling matter for the book. First, the 1847
loan was sold \emph{almost entirely in the United States}, not abroad. Only when
the domestic market sagged in the slump of 1848 did Corcoran carry the
\emph{1848} loan to Europe---and only after Secretary Walker advanced him \$5
million of Treasury funds to support the price and formally made him a government
agent \citep[54--55]{cohen1971business}. Second, the London placement was
\emph{not} chiefly through Baring: Baring Brothers, George Peabody \& Co., and
the Morgans initially told Corcoran the bonds could not be sold abroad and tried
to dissuade him. He went anyway, and the London tranche---``more than \$3
million,'' roughly a fifth of the \$16 million loan, plus some placed in
Paris---was syndicated across \emph{six} houses: Baring Brothers, George Peabody
\& Co., Overend, Gurney \& Co., Denison \& Co., Samuel Jones Loyd, and James
Morrison \citep[54--58]{cohen1971business}. The London sale itself was essentially
break-even once exchange was reckoned; the real profit came at \emph{home}, where
the news that the bonds had been placed in London restored confidence and let
Corcoran sell his remaining stock at a premium \citep{goldstein1998capital}. Only
afterward did Baring become a cooperative partner---``the demand for United
States stock runs away with all we have or wish to sell,'' it wrote him in 1849
\citep[57]{cohen1971business}.

\subsection{The London and Frankfurt houses}

Behind the antebellum Treasury stood the great transatlantic houses. \emph{Baring
Brothers} had been London agent for both Banks of the United States and, after a
brief interval in which the government account passed to Rothschild, had it
restored in 1843; Baring remained the pre-eminent European agent for American
government and state securities down to the Civil War \citep{hidy1949house}. Its
rival, N.\,M. Rothschild \& Sons, was represented in New York by \emph{August
Belmont}, who---stranded in New York by the Panic of 1837, the Rothschilds' prior
agent having failed---founded August Belmont \& Co. at 78 Wall Street and became
their permanent American agent \citep{ferguson1999house, black1981king,
sexton2005debtor}. The antebellum pattern was thus a domestic wholesaler
(Corcoran \& Riggs) buying whole loans, with the London houses a secondary outlet
resorted to only when the home market could not absorb the paper.

\section{Jay Cooke and the Financing of the Civil War, 1861--1865}

\subsection{The National Loan and Jay Cooke's retail machine}

The Civil War's borrowing needs dwarfed anything before. After the early
``National Loan'' and the resort to legal-tender greenbacks, Treasury Secretary
Salmon P. Chase turned to \emph{Jay Cooke \& Co.} of Philadelphia as special
agent. Cooke's innovation was \emph{mass retail distribution}: through a network
of roughly $2{,}500$ sub-agents, a saturation newspaper-advertising campaign,
and an appeal to patriotism, he sold government bonds directly to hundreds of
thousands of ordinary citizens \citep{larson1936jaycooke, oberholtzer1907jaycooke}.
Cooke marketed the ``five-twenty'' bonds authorized by the Act of 25 February
1862 (bonds redeemable in $5$ and payable in $20$ years, bearing $6$ percent),
selling the authorized \$500 million plus about \$11 million more; then, in
1865, he distributed roughly \$830 million of ``seven-thirty'' Treasury notes
in three series to carry the Union through the war's final months
\citep[128--173]{larson1936jaycooke}.

\subsection{How much went abroad?}

Unlike the antebellum loans, the Civil War debt was placed overwhelmingly at
\emph{home}: European investors, wary of Union prospects and of the greenback,
took little of the debt at issue. Some five-twenties later migrated to Europe,
especially to Germany and the Netherlands, as Northern victory became likely and
as the bonds' gold-interest feature attracted foreign buyers, but the wartime
distribution was a domestic achievement---Cooke had, in effect, sold the war to
the American middle class \citep{sexton2005debtor, hammond1970sovereignty}. The
large-scale migration of the debt to Europe came only \emph{after} the war: by
the early 1870s foreign investors held perhaps a third of the outstanding
bonds---Cooke's own 1871 circular claimed some \$600 million of five-twenties
were held abroad---a peak from which, as the next section shows, European
holdings then \emph{receded} through the refunding decade.

\section{From Cooke to Morgan: The Refunding Syndicates, 1865--1879}

\subsection{The Funding Acts and the breaking of Cooke's privilege}

The Funding Act of 14 July 1870, as amended 20 January 1871, authorized the
conversion of the high-interest war debt into new coin bonds in three classes:
\$500 million at $5$ percent (redeemable after ten years), \$300 million at
$4\tfrac{1}{2}$ percent (after fifteen), and \$1{,}000 million at $4$ percent
(after thirty), the aggregate capped at \$1.5 billion
\citep[477--478]{morgan1885}.\footnote{The original act of
July 1870 set the $5$ percent class at \$200 million; the January 1871 amendment
raised it to \$500 million without increasing the aggregate.} Jay Cooke \& Co.
were the Treasury's agents for the first \$200 million of $5$ percent bonds in
1871 \citep[315--330]{larson1936jaycooke}. Cooke's circular of August 1871
announced that of the remaining \$130
million, \$80 million was ``reserved for the Foreign market'' (with \$20 million
said to be subscribed in London) and \$50 million for domestic national banks---
but these were \emph{allotment and promotional} figures, not results. The
Franco-Prussian War had drained European funds, and the realized foreign uptake
was far smaller (contemporary accounts suggest only some \$15--25 million), the
bulk being taken at home \citep{dewey1903financial}. The disappointing foreign
result is precisely why the Treasury next forced Cooke into a European syndicate.
In the \$300 million $5$ percent funded loan of 1873---whose London prospectus of
31 January 1873 issued from Rothschild's New Court and named Baring Brothers,
Jay Cooke, McCulloch \& Co., J.\,S. Morgan \& Co., Morton, Rose \& Co., and
N.\,M. Rothschild \& Sons as receiving houses---the Treasury \emph{forced Cooke
to share} the flotation, ending his privileged position even before his firm fell
\citep[36]{chernow1990house}. That flotation largely failed: the post-war
European glut and then the Panic of 1873 meant little of the \$300 million was
placed, and no bonds were sold below $5$ percent until 1877
\citep[356]{dewey1903financial}.\footnote{On the failed 1873 syndication see
\citet[394--402]{larson1936jaycooke}.} On 18 September 1873,
overcommitted to the Northern Pacific Railroad, Jay Cooke \& Co. suspended and
declared bankruptcy, touching off the Panic of 1873 and the ensuing ``Long
Depression'' \citep[409--416]{larson1936jaycooke}.

\subsection{Drexel, Morgan \& Co. and the syndicates of 1877 and 1879}

Cooke's fall opened the field to \emph{Drexel, Morgan \& Co.}, formed 1 July
1871 by the union of Anthony J. Drexel of Philadelphia and the younger
J.\,Pierpont Morgan, and backed in London by Morgan's father's house, \emph{J.\,S.
Morgan \& Co.} (successor to George Peabody \& Co.)
\citep[22--29]{carosso1970investment}. Its transatlantic reach made it the
indispensable American member of every subsequent refunding group, ``usurp[ing]
[Cooke's] exalted place in government finance'' \citep[35--37]{chernow1990house}.
The refunding operations Drexel, Morgan and its allies carried through to 1879
ultimately retired or exchanged some \$1.4 billion of the debt
\citep[176--199]{carosso1987morgans}. Under the Refunding Act, Secretary
John Sherman's syndicate had been selling $4\tfrac{1}{2}$ percent bonds under a
contract of August 1876; on 9 June 1877 Sherman contracted with the same
associates for $4$ percent consols, and when subscriptions opened on 16 June
1877 they reached \$75{,}496{,}550 within thirty days, of which \$50 million was
applied to redeem an equal amount of $6$ percent bonds; some \$128.7 million more
of the $4$ percents were sold in 1878 \citep{treasury1877}. The 1877 report names
no firms---it speaks only of ``the associated contractors''---but they were the
standing syndicate documented in the parallel 1875 and 1879 contracts. That
best-documented group is the syndicate of the
refunding contract of 21 January 1879, whose members are recited verbatim in the
Supreme Court's opinion in \emph{Morgan v. United States} (1885)
\citep{morgan1885}:
\begin{itemize}
\item \emph{August Belmont \& Co.}, New York, for N.\,M. Rothschild \& Sons, London;
\item \emph{Drexel, Morgan \& Co.}, New York, for J.\,S. Morgan \& Co., London;
\item \emph{J.\,\& W. Seligman \& Co.}, New York, for Seligman Brothers, London;
\item \emph{Morton, Bliss \& Co.}, New York, for Morton, Rose \& Co., London.
\end{itemize}
Each New York house acted expressly ``on behalf of'' a London partner, and the
Treasury even established a delivery agency in London---the syndicates were a
genuinely transatlantic distribution machine. Yet it would be a mistake to infer
that the refunding \emph{placed} the American debt in Europe. On the contrary,
the documentary record shows Europe was, on net, a \emph{seller}: foreign
holdings of the federal debt fell from roughly \$500--600 million in the early
1870s to \emph{under \$200 million} by early 1879 (Garfield, in his
specie-resumption address of January 1879, put the figure at ``less than two
hundred millions''). The $4$ percent refundings of 1877--1879 were absorbed
overwhelmingly by domestic ``popular'' subscription, and much of the gold that
financed resumption came from Europe \emph{returning} American bonds. The London
syndicates were the indispensable \emph{mechanism} and took meaningful blocks,
but the \emph{stock} of United States debt held abroad shrank over the decade
\citep{morgan1885, dewey1903financial}. The allied domestic commercial
banks---above all the First National Bank of New York under George F. Baker---
supplied mobile funds, though the tight ``money-trust'' alliance later
caricatured by the Pujo Committee belongs to the 1880s--1900s rather than the
1870s.

\section{The Gold Crises and the End of the Syndicate Era, 1879--1900}

The resumption of specie payments (1879) and the silver agitation of the 1890s
made the \emph{Treasury's gold reserve}, not fresh war borrowing, the occasion
for the era's most famous underwriting. In the crisis of 1895, with the gold
reserve draining toward exhaustion, the government contracted (8 February 1895)
with a syndicate led by \emph{J.\,P. Morgan} and \emph{August Belmont} (again for
the Rothschilds) to buy \$62{,}315{,}400 of $30$-year $4$ percent coin bonds; in
return the syndicate delivered gold worth about \$65 million---roughly $3.5$
million ounces---undertaking that at least half would be drawn from \emph{Europe}
so as to replenish rather than merely recirculate the domestic stock, and
pledging to protect the Treasury against further withdrawals for six months
\citep{cleveland1895message, ferguson1999house}. The issue was split evenly
between the American and European groups, each taking half---\$31{,}157{,}750
apiece \citep[339]{carosso1987morgans}. Because $4$ percent bonds were sold at a
premium for gold, President Cleveland reckoned the effective rate at only
$3\tfrac{3}{4}$ percent. The deal was politically explosive---a Populist symbol
of Wall Street's grip on the Treasury---and its profit to the syndicate, which
Morgan pointedly declined to disclose to the Senate, remains disputed to this day
\citep[234]{noyes1909forty}. On the operating detail see
\citet[319--340]{carosso1987morgans} and \citet[73--78]{chernow1990house}.
Partly in reaction, when the Treasury needed money again in 1896 it sold \$100
million of the same $4$ percent bonds by \emph{open public bidding} rather than
secret contract (though a Morgan syndicate still bid and won a large block)
\citep[256]{noyes1909forty}.
And the Spanish-American War of 1898 was financed by a deliberately \emph{popular}
loan---\$200 million of $3$ percent ``ten-twenty'' bonds (dated 1 August 1898),
sold in denominations as low as \$20 with preference to small subscribers. It
drew 320{,}226 subscriptions totaling more than \$1.5 billion---about
seven-and-a-half times the amount offered---foreshadowing the mass finance of
1917 \citep{noyes1909forty, dewey1903financial}. Across 1865--1914, as the debt
shrank and American capital markets deepened, the European-held share of the
federal debt trended downward, even as European investment poured into American
\emph{railroads} and \emph{state} securities.

\section{The Great Reversal: World War I and the Liberty Loans, 1914--1919}

\subsection{Morgan as the Allies' banker}

The First World War inverted the century-old relationship. Designated Britain's
sole purchasing agent in the United States in January 1915, \emph{J.\,P. Morgan
\& Co.} served as the sole purchasing and fiscal agent for both Britain and
France, buying war supplies on their account \citep[187--192]{chernow1990house}
and, in October 1915, organizing the \emph{Anglo-French Loan}---\$500 million of
five-year $5$ percent bonds, issued at $98$ and marketed to \emph{American}
investors through a syndicate of $61$ underwriters and some $1{,}570$ financial
institutions \citep[198--202]{chernow1990house}. The loan was fully underwritten
but poorly absorbed by a public in which German- and Irish-American sentiment ran
against the Allies; the bonds soon traded below par \citep{burk1985sinews}. Still, the
symbolism was decisive: for the first time a great European war was being financed
in New York rather than the other way round. Across the war the United States
swung from a net international debtor of roughly \$2.2 billion in 1914 to a net
creditor of about \$6.4 billion by 1919---and that figure excludes some \$10
billion of inter-Allied \emph{government} war loans, which made the country the
world's paramount creditor.

\subsection{The Liberty Loans: the government sells directly to the citizenry}

Once the United States entered the war, Treasury Secretary William G. McAdoo
rejected the syndicate model entirely. The four \emph{Liberty Loans} of
1917--1918 and the \emph{Victory Loan} of 1919 (Table~\ref{tab:liberty}) were
sold \emph{directly to the public}, with the newly created \emph{Federal Reserve
Banks} acting as the Treasury's fiscal agents---each Reserve Bank a central agency
in its district for organizing the drive, receiving subscriptions, and delivering
bonds---together with vast volunteer ``Liberty Loan committees,'' banks,
employers, and a propaganda apparatus of posters, rallies, and celebrity drives.
The subscriber base swelled from $4$ million on the First Loan to nearly $23$
million on the Fourth. The five drives, every one oversubscribed, together had
about \$21.4 billion accepted \citep{gilbert1970american,
kangrockoff2015capitalizing, mcadoo1931crowded, fedhistory_libertybonds}. The
private underwriter, central to federal borrowing since the Revolution, had been
displaced by the Federal Reserve and the mass market. The debt was held
overwhelmingly at home; the era of European ``wholesaling'' of the United States
government's bonds was over.

\begin{table}[htbp]
\centering
\caption{The Liberty Loans and the Victory Loan, 1917--1919 (dollar amounts in
billions)}
\label{tab:liberty}
\small
\begin{tabular}{l l c r r r}
\toprule
Loan & Dated & Coupon & Offered & Accepted & Subscribers \\
\midrule
First Liberty  & June 15, 1917 & $3.50\%$          & \$2.00 & \$2.00 & 4.0 M \\
Second Liberty & Nov.\ 15, 1917 & $4.00\%$          & \$3.00 & \$3.81 & 9.4 M \\
Third Liberty  & May 9, 1918   & $4.25\%$          & \$3.00 & \$4.18 & 18.4 M \\
Fourth Liberty & Oct.\ 24, 1918 & $4.25\%$          & \$6.00 & \$6.96 & 22.8 M \\
Victory (Fifth) & May 1919     & $4.75\%/3.75\%$   & \$4.50 & \$4.50 & 11.8 M \\
\midrule
\textbf{Total} & & & & \textbf{\$21.4} & \\
\bottomrule
\end{tabular}

\smallskip
\footnotesize Source: \citet{kangrockoff2015capitalizing}, Table 1. ``Accepted''
is the amount the Treasury allotted; every loan was oversubscribed. The Victory
notes were issued in two series, a $4\tfrac{3}{4}$ percent tax-exempt and a
$3\tfrac{3}{4}$ percent taxable.
\end{table}

\section{Peacetime Debt Management and the WWII Coda, 1919--1945}

Through the 1920s the Treasury managed and refunded its large war debt through
regular offerings sold via the Federal Reserve Banks and a growing market of
\emph{government-securities dealers} and commercial banks---routine, competitive
distribution rather than privileged underwriting \citep{studenski1963financial,
myers1970financial}. Foreign holdings of the federal debt were by now negligible;
the United States was the world's leading creditor. The Second World War repeated
and enlarged the WWI template: a continuous \emph{War Finance} program
(1941--1945) sold \$185.7 billion of securities, of which the retail
\emph{Series~E} ``war bonds'' sold to individuals totaled \$33.7 billion (all
savings bonds, about \$54 billion), purchased by some $85$ million Americans,
with the Federal Reserve Banks again acting as fiscal agents and the debt
absorbed essentially entirely at home \citep{treasury1945, treasurysavings1984,
bickley2010warbonds}.

\section{Synthesis: The Marketers and the European Share}

Table~\ref{tab:marketers} summarizes the chief marketers of federal debt by
regime; Table~\ref{tab:foreign} collects the estimated European-held fraction of
the federal debt at benchmark dates. Two long movements stand out. First, the
\emph{intermediary} evolved from foreign sovereigns and bankers, to a national
bank as fiscal agent, to domestic private wholesalers, to transatlantic
investment-banking syndicates, and finally to the Federal Reserve selling
directly to the citizenry. Second, the \emph{European share} did not move in one
direction but oscillated: rising from perhaps a fifth in 1789 to a peak of
$53$--$56$ percent in 1803, receding as the debt was paid off before 1836,
reviving to about a third when Civil War five-twenties migrated abroad in the late
1860s, then \emph{receding again} through the 1870s (when, contrary to a common
impression, Europe was a net seller and the refunding syndicates were a
distribution channel rather than a net absorber), before collapsing toward zero
after 1917 as the United States turned from the world's largest sovereign debtor
into its largest creditor.

\begin{table}[htbp]
\centering
\caption{Chief marketers / wholesalers of U.S. federal debt, by regime}
\label{tab:marketers}
\small
\begin{tabular}{p{2.5cm} p{6.5cm} p{4.5cm}}
\toprule
Era & Principal marketer(s) / underwriter(s) & Mechanism \\
\midrule
1776--1790 & State \& Continental loan offices; Robert Morris \& Bank of North America; France; Amsterdam consortium (Willink, Van Staphorst, De la Lande \& Fynje); Haym Salomon (broker) & Loan-office certificates; foreign sovereign loans; Dutch bond flotations \\
1790--1811 & First Bank of the United States (fiscal agent); Dutch \& British investors; Baring \& Hope (Louisiana stock) & National bank as agent; European placement \\
1812--1836 & Girard, Astor \& Parish (1813 syndicate); Second Bank of the U.S. (fiscal agent) & First private war-loan syndicate; national-bank agency \\
1836--1861 & Corcoran \& Riggs (wholesaler); London houses---Baring, Peabody, Overend Gurney, and others---as a secondary outlet; August Belmont (Rothschild agent) & Wholesale purchase, chiefly domestic; London resorted to only when the home market sagged \\
1861--1865 & Jay Cooke \& Co. & Mass retail to U.S. citizens via sub-agents \\
1865--1879 & Jay Cooke (to 1873); then Drexel, Morgan \& Co. / J.\,S. Morgan; Belmont/Rothschild; Seligman; Morton & Refunding syndicates, London-backed \\
1879--1900 & J.\,P. Morgan--Belmont syndicate (1895 gold rescue); direct public sale (1896, 1898) & Syndicate; then popular subscription \\
1917--1919 & Federal Reserve Banks \& Liberty Loan committees (direct); J.\,P. Morgan (agent for the Allies) & Mass public subscription; U.S. as creditor \\
1919--1945 & Federal Reserve Banks; securities dealers; direct public (war bonds) & Competitive distribution; retail savings bonds \\
\bottomrule
\end{tabular}
\end{table}

\begin{table}[htbp]
\centering
\caption{Estimated foreign (mainly European) share of the U.S. federal debt}
\label{tab:foreign}
\small
\begin{tabular}{c c p{7.5cm}}
\toprule
Year & Foreign share & Notes and source \\
\midrule
1789 & $\approx 22$--$29\%$ & Foreign-currency loans (France, Dutch) $\approx$\$11.7M of $\approx$\$54M; higher figure adds domestic debt bought abroad \citep{wilkins1989history, bayley1881national} \\
1803 & $\approx 53$--$56\%$ & Peak; \$43.4M of \$81.4M federal debt. Louisiana sixes taken up essentially entirely abroad \citep{syllawilsonwright2006, wilkins1989history} \\
1818 & $\approx 26\%$ & After partial pay-down \citep{wilkins1989history} \\
1835 & $0\%$ & Federal debt extinguished \citep{lane2014nation} \\
1853 & $\approx 46\%$ & Debt held abroad after renewed borrowing \citep{wilkins1989history} \\
1861--1865 & small at issue & Civil War debt placed overwhelmingly at home; five-twenties migrated to Europe only after $\sim$1863 \citep{sexton2005debtor} \\
$\sim$1870 & $\approx \tfrac{1}{3}$ & Postwar peak: $\sim$\$500--600M of $\sim$\$1.6--1.8B held abroad \citep{morgan1885, dewey1903financial} \\
1879 & $< \$200$M & Europe a \emph{net seller} through the refunding decade; foreign stock \emph{fell} even as London syndicates distributed the bonds \citep{dewey1903financial} \\
1895 & (gold, not debt) & $\geq$ half the \$62.3M bond deal's $\sim$\$65M of gold drawn from Europe \citep{cleveland1895message, chernow1990house} \\
1917-- & $\to 0$ & U.S. becomes net creditor; debt held domestically \citep{gilbert1970american} \\
\bottomrule
\end{tabular}
\end{table}

\paragraph{Caveats for the book.} The benchmark foreign-share figures rest on
Blodget's (1806) and Bayley's Treasury estimates as reconciled by
\citet{syllawilsonwright2006} and \citet{wilkins1989history}; the 1803 figure is
best stated as a $53$--$56$ percent range because the underlying totals differ.
No clean benchmark exists for c.\,1795 or c.\,1815---the record has gaps around
those years, and any single number there is interpolation. Page references to the
Carosso, Chernow, and Larson monographs have been verified against the
first-edition pagination via Google Books' search-within-volume index (see the
notes in the \texttt{.bib}); the Ferguson and Catterall page references still
warrant a physical-volume check. Two cautions: the ``\$150{,}000 underwriters'
fee'' often attached to the 1873 Cooke syndicate does \emph{not} appear in
Chernow, Carosso, or Larson and should not be printed without an identified
source; and the exact van Winter loan-by-loan Dutch figures (Table~\ref{tab:dutch}
follows Bayley) could not be read directly and are corroborated only in
aggregate. The syndicate rosters, loan schedules, and dollar figures in the text
were otherwise, where marked, checked against the \emph{Morgan v. United States}
opinion, Bayley's census monograph, and the Treasury Annual Reports.

\bibliographystyle{aer}
\bibliography{bond_underwriters}

\end{document}
